\chapter{Important C Data Structures}
\label{ann:qemu_c_structures}

\section{QOM API Core Elements} 

\textbf{ObjectClass:} Analogous to a class in any OOP language, the C \\ 
\texttt{ObjectClass} structure represents the metadata that describe a type and that is 
shared by all instances of that type. It holds\cite{Qemu_QOM_API}:

\lstinputlisting[
  language=C,
  caption={\texttt{ObjectClass}},
  label={code:ObjectClass}
]{Resources/Code/qom/ObjectClass.c}

\begin{itemize}
    \item \textbf{Type information}: Unique identifier for runtime type checking
    \item \textbf{Virtual method pointers}: Functions that can be overridden by derived types
    \item \textbf{Interface implementations}: Lists of interfaces this type provides.
    \item \textbf{Shared class properties}: Default values and metadata for all instances
\end{itemize}

\textbf{Object:} Every object in QEMU inherits from the C \texttt{Object} structure, which 
is the fundamental component of the entire type system. Each individual object 
is an \textbf{object instance} with its own state, which can be instantiated, configured, 
introspected, and destroyed. Each \texttt{Object} instance encapsulates:

\lstinputlisting[
  language=C,
  caption={\texttt{Object}},
  label={code:Object}
]{Resources/Code/qom/Object.c}

\begin{itemize}
    \item \textbf{Type identification} via a class pointer
    \item \textbf{Reference counting} for memory management
    \item \textbf{Dynamic properties} storage and access
    \item \textbf{Composition tree relationships} (parent pointers) enabling hierarchical 
                                        organization
\end{itemize}

\textbf{TypeInfo:} The C \texttt{TypeInfo} structure is the complete blueprint for registering 
a new QOM type at QEMU startup. This structure tells QEMU exactly how to create, initialize, 
and manage instances of the QOM type.

\lstinputlisting[
  language=C,
  caption={\texttt{TypeInfo}},
  label={code:TypeInfo}
]{Resources/Code/qom/TypeInfo.c}

\begin{itemize}
    \item Type name and parent class (e.g., parent=\texttt{"sysbus-device"})
    \item Instance and class sizes (memory layout)
    \item Constructor (\texttt{instance\_init}), post-constructor (\texttt{instance\_post\_init}), and 
            destructor (\texttt{instance\_finalize}) callbacks
    \item Class initialization callback (\texttt{class\_init}) where virtual methods are assigned
\end{itemize}

To exemplify the use of the \texttt{TypeInfo} structure: The \texttt{STM32I2CClass} is registered via \texttt{TypeInfo} with 
parent=\texttt{'sysbus-device'}, \texttt{instance\_size=sizeof(STM32I2CState)}, and 
\texttt{class\_init} \texttt{=stm32\_i2c\_class\_init}. This registration enables QEMU to instantiate the I$^2$C 
controller and configure its base address, interrupt lines, and master/slave modes at 
runtime.

\textbf{ObjectProperty:} Each object's configurable properties—e.g., base addresses, 
enable bits, interrupt vectors—are registered via \texttt{ObjectProperty} descriptors. These 
define getters/setters, type metadata, and default values, enabling firmware 
configuration without recompilation.

\section{QEMU Memory}

\lstinputlisting[
  language=C,
  caption={\texttt{MemoryRegion}},
  label={code:MemoryRegion}
]{Resources/Code/memory/memoryregion.c}